\section{Results}

\subsection{Two development passes become two unseen failures}

Both new equation-discovery mechanisms pass the development screen. Round 1
has a median development improvement of \RouteOneDevelopment{}, and Round 2
has \RouteTwoDevelopment{}. If evaluation stopped at development performance,
either round could be narrated as promising. The frozen unseen panels reverse
that conclusion. Round 1 obtains a system-level 95\% interval of
[\RouteOneCILower{}, \RouteOneCIUpper{}], and Round 2 obtains
[\RouteTwoCILower{}, \RouteTwoCIUpper{}]. Both intervals cross zero, so both
rounds are negative under the preregistered decision.

\begin{figure}[t]
  \centering
  \includegraphics[width=\columnwidth]{figures/route-a-results.pdf}
  \Description{Two horizontal confidence intervals cross a vertical zero
  boundary, so neither equation-discovery round passes its unseen gate.}
  \caption{\routea{} unseen results. Points show median failure-aware relative
  improvement and bars show frozen system-level bootstrap intervals. The
  dashed line is the acceptance boundary; both lower endpoints are below
  zero.}
  \label{fig:route-a}
\end{figure}

\begin{table}[t]
\centering
\caption{Development and unseen decisions for the two parent-linked mechanism
rounds. Full mechanism identifiers are retained for auditability.}
\label{tab:route-a}
\resizebox{\columnwidth}{!}{\input{tables/route-a-results}}
\end{table}

This result matters in two ways. Scientifically, the tested smoothing and
support-selection mechanisms do not establish robust improvement over the
strong baseline on their held-out systems. Operationally, the campaign does
not reuse a panel, tune after reveal, or lower the confidence threshold.
Round 1's negative hash becomes the parent of Round 2, and the second negative
decision triggers the preregistered systems route. The paper therefore
contains actual failed evidence that changed the next research object.

\subsection{Main controller comparison}

Figure~\ref{fig:main-results} reports task success for the three main
controllers. The \oneshot{} controller succeeds on \OneShotSuccess{} of
cells, \executeonce{} on \ExecuteOnceSuccess{}, and the \systemloop{} on
\FullLoopSuccess{}. Relative to \executeonce{}, the paired mean gain is
\PairedGain{}. The \BootstrapResamples{}-resample 95\% interval is
[\PairedCILower{}, \PairedCIUpper{}], whose lower endpoint is above zero.
All 30 paired differences are included in the frozen result and recomputed by
the standalone reproduction script.

\begin{figure}[t]
  \centering
  \includegraphics[width=\columnwidth]{figures/main-results.pdf}
  \Description{Bar chart with success rates 0.20 for one-shot, 0.50 for
  execute-once, and 1.00 for the full loop.}
  \caption{Task success in the preregistered main comparison. Each rate
  aggregates ten tasks and three seeds.}
  \label{fig:main-results}
\end{figure}

The \systemloop{} starts with 24 negative cells and recovers 15, for a recovery
rate of \FullLoopRecovery{}. Execute-once repairs planning-visible faults but
does not recover a runtime failure, so its measured recovery rate is zero.
The full-loop scientific and reproduction hashes agree for every cell,
yielding exact reproduction \FullLoopReproduction{}. Its reports contain
\FullLoopUnsupported{} unsupported claims, and the campaign records
\HumanInterventions{} research-decision human interventions.

\subsection{Component ablations}

Figure~\ref{fig:ablations} and Table~\ref{tab:all-modes} show the four
preregistered ablations. Removing \vault{} memory reduces success from
\FullLoopSuccess{} to \NoVaultSuccess{} because memory-dependent evidence-map
and checkpoint repairs remain blocked. Removing failure feedback produces
\NoFeedbackSuccess{}, matching execute-once because the controller cannot
act on runtime diagnosis. Removing preregistration produces
\NoPreregSuccess{} under a metric that requires the frozen contract.
Removing the evidence gate retains \NoGateSuccess{} apparent success but
allows \NoGateUnsupported{} unsupported claims. This last contrast shows why
task completion alone is insufficient as a research-system endpoint.

\begin{figure}[t]
  \centering
  \includegraphics[width=\columnwidth]{figures/ablation-results.pdf}
  \Description{Bar chart comparing the full-loop success rate of 1.00 with
  four ablations: no Vault 0.70, no feedback 0.50, no preregistration 0.00,
  and no evidence gate 0.80.}
  \caption{Task success for the complete loop and four frozen ablations.
  The no-evidence-gate bar must be read with its six unsupported claims.}
  \label{fig:ablations}
\end{figure}

\begin{table*}[t]
\centering
\caption{All controller modes. Recovery is conditional on an initially failed
cell. Reproduction compares the primary and exact-rerun scientific hashes.
Unsupported is a count over 30 cells per mode.}
\label{tab:all-modes}
\small
\input{tables/mode-results}
\end{table*}

\subsection{Local execution and evidence completeness}

The main controller summaries cause \LocalRequests{} requests to the local
model. \LocalFallbacks{} request reaches the timeout boundary and uses the
predeclared deterministic fallback. The fallback is visible in the policy
evidence and does not affect numerical adjudication. Scientific computation
and model inference remain local; the recorded external cost is zero.

Every matrix cell has a result JSON, a reproduction JSON, a Markdown report,
and hashes of its source and policy evidence. The aggregate package contains
the frozen preregistration, fault policy, evaluator hash, all \CellCount{}
cells, a mode table, figures, failure analysis, loop report, and contribution
decision. The final paper builder additionally checks bibliography resolution,
claim-evidence links, LaTeX references, vector-figure compilation, layout
overflow, and a clean-directory rebuild.
